La formación en robótica e inteligencia artificial en el ámbito universitario enfrenta un desafío persistente: la distancia entre los fundamentos teóricos impartidos en el aula y la experimentación práctica con sistemas reales que integren percepción, control y toma de decisiones. En la Universidad Nacional de Colombia, el Laboratorio de Sistemas Inteligentes Robóticos (LabSIR) cuenta con manipuladores PhantomX Pincher como plataforma principal para la enseñanza de la robótica. Estos robots, basados en servomotores Dynamixel y una arquitectura modular de código abierto, han demostrado ser herramientas valiosas para el aprendizaje de cinemática, dinámica y control de manipuladores. No obstante, su uso se ha limitado predominantemente a ejercicios de posicionamiento y trayectorias preprogramadas, sin incorporar de manera sistemática capacidades de percepción artificial ni tareas autónomas de clasificación de objetos.

En el mercado internacional, empresas como Elephant Robotics han desarrollado soluciones integradas ---como el myCobot AI Kit--- que combinan manipulación robótica, visión por computador y módulos de inteligencia artificial en un entorno compacto y didáctico. Estos kits permiten a los estudiantes experimentar con pipelines completos de percepción-acción, desde la captura y procesamiento de imágenes hasta la ejecución de movimientos de agarre y clasificación. Sin embargo, estas soluciones comerciales están diseñadas para plataformas propietarias y no son compatibles con el PhantomX Pincher, lo cual impide su adopción directa en el laboratorio. Además, su costo elevado y la dependencia de ecosistemas cerrados limitan las posibilidades de personalización y adaptación a las necesidades pedagógicas específicas del entorno universitario local.

Por otra parte, la transición de ROS (Robot Operating System) hacia ROS2 ha introducido mejoras significativas en términos de comunicación en tiempo real, seguridad y escalabilidad, posicionándose como el estándar emergente para el desarrollo de aplicaciones robóticas. Paralelamente, la Raspberry Pi 5 ofrece capacidades de procesamiento embebido suficientes para ejecutar tareas de visión por computador con bibliotecas como OpenCV, a un costo accesible para contextos académicos. A pesar de la disponibilidad de estas tecnologías de código abierto, no se ha documentado ampliamente en la literatura un kit educativo que integre de forma sistemática el PhantomX Pincher con ROS2, Raspberry Pi y un sistema de visión artificial para tareas de clasificación de objetos. Esta ausencia representa una brecha tanto en la oferta de herramientas didácticas disponibles como en el aprovechamiento del potencial de la infraestructura existente en el laboratorio.

La carencia de un entorno estructurado que articule hardware, software y guías pedagógicas tiene consecuencias concretas en la formación de los estudiantes. Sin un kit integrado, los alumnos deben invertir un tiempo desproporcionado en la configuración, compatibilidad y conexión de componentes heterogéneos, lo cual reduce el tiempo efectivo dedicado al aprendizaje de conceptos fundamentales como la visión por computador, la planificación de movimiento y el control autónomo. Asimismo, la falta de documentación estandarizada y de un diseño replicable dificulta la escalabilidad de las prácticas de laboratorio y genera una curva de aprendizaje innecesariamente pronunciada.

Ante este panorama, surge la siguiente pregunta de investigación: \textbf{¿Cómo diseñar un kit educativo de clasificación de objetos con visión por computador para el robot PhantomX Pincher que, utilizando ROS2 y Raspberry Pi 5, integre de manera modular y replicable los subsistemas de percepción, control y manipulación, y que resulte adecuado para su implementación en entornos académicos de enseñanza de la robótica?}

Este trabajo busca abordar dicha problemática mediante el diseño conceptual de un kit que contemple tanto los aspectos mecánicos ---plataforma base, zonas de recolección y depósito, soporte de cámara, sistema de succión y base del manipulador--- como los aspectos de software ---arquitectura de nodos ROS2, procesamiento visual con OpenCV y planificación de movimiento con MoveIt2---. El diseño se orienta hacia la producción de siete unidades dentro de restricciones presupuestarias definidas, empleando tecnologías abiertas y procesos de fabricación accesibles como la impresión 3D (FDM), con el objetivo de sentar las bases documentales y técnicas para su futura implementación en los cursos de robótica de la universidad.